% =====================================================================
%  Back cover — Is the Riemann Hypothesis Solved?
%  Typeset standalone at the 6 x 9 in trim, full bleed.
%  Vertical positions are measured in inches from the foot of the page.
% =====================================================================
\documentclass{article}
\usepackage[paperwidth=6in,paperheight=9in,margin=0pt]{geometry}
\usepackage{fontspec}
\usepackage{xcolor}
\usepackage{graphicx}
\usepackage{tikz}
\usepackage{microtype}

\setmainfont{Outfit}[
  UprightFont={Outfit Regular}, BoldFont={Outfit Bold},
  ItalicFont={Outfit Regular}, ItalicFeatures={FakeSlant=0.20},
  BoldItalicFont={Outfit Bold}, BoldItalicFeatures={FakeSlant=0.20},
  Ligatures={TeX}]
% The cover sets most of its type in \sffamily. Without an explicit
% sans face that is Latin Modern Sans, so the cover would carry two
% typefaces where the book carries one.
\setsansfont{Outfit}[
  UprightFont={Outfit Regular}, BoldFont={Outfit Bold},
  ItalicFont={Outfit Regular}, ItalicFeatures={FakeSlant=0.20},
  BoldItalicFont={Outfit Bold}, BoldItalicFeatures={FakeSlant=0.20},
  Ligatures={TeX}]
\setmonofont{DejaVu Sans Mono}[
  UprightFont={DejaVu Sans Mono}, BoldFont={DejaVu Sans Mono Bold},
  Ligatures={TeX}]

\definecolor{ink}{HTML}{0C1016}
\definecolor{accent}{HTML}{8FD6C0}
\definecolor{txt}{HTML}{E6EBEF}
\definecolor{dim}{HTML}{8B9BA7}

\pagestyle{empty}
\setlength{\parindent}{0pt}

\begin{document}
\begin{tikzpicture}[remember picture,overlay]
  % ---- the plate, full bleed (a different crop from the front) ----------
  \node[anchor=south west,inner sep=0] at (current page.south west)
    {\includegraphics[width=\paperwidth,height=\paperheight]{../figures/cover_back_art.jpg}};



  % ---- title strip -------------------------------------------------------
  \node[anchor=south west,inner sep=0,text=accent] at ([xshift=0.60in,yshift=8.26in]current page.south west)
    {\fontsize{9.2}{11}\selectfont\sffamily\bfseries
     {\addfontfeatures{LetterSpace=15}IS THE RIEMANN HYPOTHESIS SOLVED?}};

  \fill[accent] ([xshift=0.60in,yshift=8.00in]current page.south west)
         rectangle ([xshift=5.48in,yshift=8.02in]current page.south west);

  \node[anchor=south west,inner sep=0,text=dim] at ([xshift=0.60in,yshift=7.66in]current page.south west)
    {\fontsize{9.8}{12}\selectfont\sffamily Gerald Enrique Nelson Mc Kenzie};

  % ---- blurb -------------------------------------------------------------
  \node[anchor=north west,inner sep=0,text=txt,align=left] at ([xshift=0.60in,yshift=7.22in]current page.south west)
    {\fontsize{8.3}{11.4}\selectfont\sffamily
     In November 2026 a mathematician published four pages claiming to have\\
     proved the Riemann Hypothesis, working with Claude, together with a Lean\\
     development of 175,187 files and 17.4 million lines. Two university teams\\
     compiled it the same day. The world's newspapers called it the greatest\\
     mathematical result of the century.\\[7pt]
     On 21 November the Lean and Mathlib communities certified it. The\\
     statement is Mathlib's own; the exported proof is accepted by Lean's\\
     kernel and by four independently written checkers.\\[7pt]
     \textbf{No mathematician has read the argument.}\\[7pt]
     This book takes the artifact seriously, reads it, and traces the proof\\
     chain from Mathlib's own definition of the hypothesis down to a proved\\
     fixed-point formula. It identifies the load-bearing inputs that nobody\\
     has audited, and sets out the author's own admissions: a computed\\
     ceiling at two thirds of the zeros, and an obstruction whose stated\\
     remedy requires a conjecture that is still open.\\[7pt]
     The finding is neither applause nor dismissal. The claim is formally\\
     certified and mathematically unaudited, and Chapter 14 gives the\\
     protocol that would settle it.\\
     \\
     The mathematics examined here is the work of Levent Alpöge. This book\\
     is the analysis and exposition of it.};

  % ---- parts -------------------------------------------------------------
  \node[anchor=south west,inner sep=0,text=accent] at ([xshift=0.60in,yshift=3.18in]current page.south west)
    {\fontsize{6.8}{8}\selectfont\sffamily\bfseries {\addfontfeatures{LetterSpace=18}INSIDE}};

  \node[anchor=north west,inner sep=0,text=txt,align=left] at ([xshift=0.60in,yshift=3.04in]current page.south west)
    {\fontsize{8.6}{11.6}\selectfont\sffamily
     I \quad The Claim\\
     II \quad The Mathematics\\
     III \quad The Method, Step by Step\\
     IV \quad The Proof in Natural Language\\
     V \quad The Verification\\
     VI \quad The State of Play\\
     VII \quad The Verdict};

  % ---- footer ------------------------------------------------------------
  \fill[dim,opacity=0.45] ([xshift=0.60in,yshift=0.92in]current page.south west)
         rectangle ([xshift=5.48in,yshift=0.93in]current page.south west);

  \node[anchor=south west,inner sep=0,text=dim] at ([xshift=0.60in,yshift=0.66in]current page.south west)
    {\fontsize{7.6}{9.6}\selectfont\sffamily
     Tico Internet LLC \quad\textperiodcentered\quad 2026 First Edition
     \quad\textperiodcentered\quad Apache License 2.0};

  \node[anchor=south west,inner sep=0,text=dim] at ([xshift=0.60in,yshift=0.46in]current page.south west)
    {\fontsize{7.6}{9.6}\selectfont\sffamily DOI 10.5281/zenodo.24408811};
\end{tikzpicture}
\null
\end{document}
